% Options for packages loaded elsewhere
\PassOptionsToPackage{unicode}{hyperref}
\PassOptionsToPackage{hyphens}{url}
\documentclass[
]{article}
\usepackage{xcolor}
\usepackage{amsmath,amssymb}
\setcounter{secnumdepth}{-\maxdimen} % remove section numbering
\usepackage{iftex}
\ifPDFTeX
  \usepackage[T1]{fontenc}
  \usepackage[utf8]{inputenc}
  \usepackage{textcomp} % provide euro and other symbols
\else % if luatex or xetex
  \usepackage{unicode-math} % this also loads fontspec
  \defaultfontfeatures{Scale=MatchLowercase}
  \defaultfontfeatures[\rmfamily]{Ligatures=TeX,Scale=1}
\fi
\usepackage{lmodern}
\ifPDFTeX\else
  % xetex/luatex font selection
\fi
% Use upquote if available, for straight quotes in verbatim environments
\IfFileExists{upquote.sty}{\usepackage{upquote}}{}
\IfFileExists{microtype.sty}{% use microtype if available
  \usepackage[]{microtype}
  \UseMicrotypeSet[protrusion]{basicmath} % disable protrusion for tt fonts
}{}
\makeatletter
\@ifundefined{KOMAClassName}{% if non-KOMA class
  \IfFileExists{parskip.sty}{%
    \usepackage{parskip}
  }{% else
    \setlength{\parindent}{0pt}
    \setlength{\parskip}{6pt plus 2pt minus 1pt}}
}{% if KOMA class
  \KOMAoptions{parskip=half}}
\makeatother
\setlength{\emergencystretch}{3em} % prevent overfull lines
\providecommand{\tightlist}{%
  \setlength{\itemsep}{0pt}\setlength{\parskip}{0pt}}
\usepackage{bookmark}
\IfFileExists{xurl.sty}{\usepackage{xurl}}{} % add URL line breaks if available
\urlstyle{same}
% fallback for those not using the hyperref driver hyperxmp:
\makeatletter
\@ifundefined{xmpquote}{\newcommand{\xmpquote}[1]{#1}}{}
\makeatother
\hypersetup{
  hidelinks,
  pdfcreator={LaTeX via pandoc}}

\author{}
\date{}

\begin{document}

\textbf{ANALISIS PENURUNAN TEKANAN DAN TRANSFER KALOR \emph{HEAT PIPE}
DI REAKTOR ANTARIKSA MENGGUNAKAN PYTHON}

\textbf{SKRIPSI}

untuk memenuhi sebagian persyaratan

untuk memperoleh derajat Sarjana

Program Studi Teknik Nuklir

Diajukan oleh

MUHAMMAD NAUFAL TAMIMI

20/456315/TK/50445

Kepada

\textbf{DEPARTEMEN TEKNIK NUKLIR DAN TEKNIK FISIKA}

\textbf{FAKULTAS TEKNIK}

\textbf{UNIVERSITAS GADJAH MADA}

\textbf{YOGYAKARTA}

\textbf{2026}

\emph{\ldots Karya ini kupersembahkan untuk keluarga dan teman-teman
yang menemaniku\ldots{}}

\section{\texorpdfstring{\textbf{BAB I}
\textbf{PENDAHULUAN}}{BAB I PENDAHULUAN}}\label{bab-i-pendahuluan}

\begin{enumerate}
\def\labelenumi{\arabic{enumi}.}
\item ~
  \subsection{\texorpdfstring{\textbf{Latar
  Belakang}}{Latar Belakang}}\label{latar-belakang}
\end{enumerate}

Eksplorasi antariksa yang semakin masif, baik untuk misi penjelajahan
antariksa dalam maupun pembangunan pangkalan di bulan dan planet lain,
membutuhkan sumber energi yang andal dan berkapasitas besar. Sumber daya
konvensional seperti panel surya memiliki keterbatasan karena
ketersediaan berselang dan kinerjanya yang menurun drastis saat berada
jauh dari matahari, sedangkan \emph{Radioisotope Thermoelectric
Generators} (RTG) yang menggunakan isotop seperti Plutonium-238 memiliki
kendala pada daya keluaran yang relatif kecil, biasanya di bawah 500 W.
Oleh karena itu, reaktor fisi nuklir antariksa, seperti proyek Kilopower
yang dikembangkan oleh NASA, hadir sebagai solusi menjanjikan untuk
mengisi kesenjangan kebutuhan daya antara 1 hingga 10 kWe (atau lebih)
untuk misi jangka panjang
\href{https://www.zotero.org/google-docs/?e7LRKG}{{[}1{]}}. Reaktor
antariksa menawarkan kepadatan daya yang tinggi, sistem yang sangat
kompak, serta kemampuan beroperasi secara terus-menerus hingga puluhan
tahun tanpa bergantung pada kondisi lingkungan luar.

Dalam merancang desain reaktor mikro antariksa, sistem pendingin dan
transfer kalor yang efisien sekaligus aman merupakan tantangan teknis
yang krusial. Pipa kalor (\emph{heat pipe}) logam alkali suhu tinggi,
seperti yang menggunakan natrium cair, telah menjadi teknologi pemindah
kalor pasif utama untuk aplikasi reaktor antariksa. Perangkat
termodinamika tertutup ini beroperasi dengan memanfaatkan kalor laten
penguapan dari fluida kerjanya untuk memindahkan jumlah energi kalor
yang signifikan dengan gradien suhu yang sangat minimal. Proses aliran
balik cairan murni digerakkan oleh tekanan kapiler yang dihasilkan di
dalam struktur sumbu (\emph{wick}) berpori, sehingga menghilangkan
kebutuhan akan pompa mekanis aktif. Pengoperasian pasif ini tidak hanya
sangat andal, tetapi juga mencegah kemungkinan kegagalan komponen
mekanik tunggal dan sangat ideal untuk beroperasi di lingkungan
gravitasi mikro di luar angkasa
\href{https://www.zotero.org/google-docs/?0U09MA}{{[}2{]}}.

Meskipun memiliki efisiensi transfer kalor yang sangat tinggi, operasi
pipa kalor sangat dibatasi oleh batasan-batasan fisik yang nonlinier.
Batasan ini meliputi batas sonik saat aliran uap tercekik (\emph{choked
flow}), batas \emph{entrainment} karena gesekan uap berkecepatan tinggi,
batas viskositas pada saat suhu awal rendah, batas pendidihan
(\emph{boiling limit}), dan batas kapiler. Secara khusus, batas kapiler
sering kali menjadi pembatas utama kapasitas perpindahan kalor pada
keadaan tunak (\emph{steady-state}); batas ini tercapai apabila tekanan
kapiler tidak lagi mampu mengatasi total penurunan tekanan
(\emph{pressure drop}) yang timbul dari aliran cairan di dalam sumbu,
gesekan uap, dan gaya gravitasi. Apabila beban kalor atau tekanan
hidrolik melebihi kemampuan desainnya, sirkulasi akan terputus dan pipa
kalor akan mengalami kekeringan pada bagian evaporator (\emph{dryout}),
yang berpotensi memicu lonjakan suhu ekstrem yang membahayakan inti
reaktor \href{https://www.zotero.org/google-docs/?QjHnZM}{{[}3,
p.~170-179{]}}.

Mengingat kompleksitas operasi tersebut, diperlukan model komputasi yang
tangguh bagi desainer termal-hidrolik untuk mengevaluasi parameter
desain pipa kalor. Secara historis, analisis pipa kalor konvensional
sering kali mengandalkan kode-kode lawas berbasis Fortran seperti HTPIPE
dari Los Alamos National Laboratory, atau kode berskala rekayasa modern
yang membutuhkan komputasi masif seperti Sockeye berbasis MOOSE. Akan
tetapi, fleksibilitas bahasa pemrograman Python memberikan peluang untuk
mengembangkan alat pemodelan komputasi yang mandiri
(\emph{self-contained}), efisien, dan komprehensif. Melalui skrip
Python, perhitungan distribusi tekanan, profil suhu uap, dan
penyelesaian batas operasi fisik (\emph{operational limits}) dapat
diprediksi secara otomatis, memungkinkan optimalisasi desain pipa kalor
yang lebih terintegrasi untuk berbagai variasi reaktor antariksa

Prototipe eksperimental Kilopower Reactor Using Stirling TechnologY
(KRUSTY) telah secara sukses mendemonstrasikan keandalan pendinginan
pipa kalor natrium pada uji daya penuh untuk reaktor antariksa. Konsep
sukses ini dapat diaplikasikan dan dioptimalkan lebih lanjut pada
desain-desain reaktor masa depan, termasuk desain konseptual Reaktor
Antariksa Indonesia atau Indonesian Space Reactor (ISR) yang dirancang
menggunakan spektrum cepat untuk menghasilkan daya 500 kWth secara
kontinu. Konfigurasi pendinginan pasif pada ISR ditujukan untuk
memastikan distribusi termal yang merata, memitigasi titik panas lokal,
serta menjamin reaktor tetap kritis dan aman dari kegagalan mekanis.

\begin{enumerate}
\def\labelenumi{\arabic{enumi}.}
\setcounter{enumi}{1}
\item ~
  \subsection{\texorpdfstring{\textbf{Perumusan
  Masalah}}{Perumusan Masalah}}\label{perumusan-masalah}
\end{enumerate}

Pipa kalor (\emph{heat pipe}) logam alkali suhu tinggi merupakan
teknologi pendingin pasif utama yang krusial untuk aplikasi reaktor
mikro terestrial maupun tenaga nuklir antariksa. Namun, operasi pipa
kalor sangat dibatasi oleh lima batasan fisik yang ketat, yaitu batas
sonik, batas viskos, batas kapiler, batas \emph{entrainment}, dan batas
pendidihan. Jika beban termal melampaui batasan ini pada suhu operasi
tertentu, pipa kalor akan gagal berfungsi, yang berujung pada lonjakan
suhu secara cepat dan fenomena kekeringan pada evaporator (dryout).
Secara historis, desainer termal-hidrolik sangat bergantung pada kode
lawas berbasis Fortran (seperti HTPIPE) atau \emph{framework} modern
yang membutuhkan komputasi masif (seperti Sockeye berbasis MOOSE) untuk
menganalisis batas-batas ini. Di sisi lain, desain konseptual Reaktor
Antariksa Indonesia (ISR) yang dirancang menghasilkan daya 500 kWth
memerlukan sistem pendingin pasif yang tangguh agar terhindar dari titik
panas lokal dan kegagalan komponen. Oleh karena itu, diperlukan sebuah
alat pemodelan komputasi yang mandiri (\emph{self-contained}), efisien,
dan presisi menggunakan bahasa pemrograman Python untuk menganalisis
batasan kinerja tersebut sebelum diaplikasikan pada desain reaktor masa
depan. Berdasarkan uraian tersebut, dapat disusun rumusan masalah
penelitian sebagai berikut:

\begin{enumerate}
\def\labelenumi{\arabic{enumi}.}
\item
  Bagaimana cara membangun dan mengembangkan model komputasi berbasis
  bahasa pemrograman Python untuk memprediksi secara akurat penurunan
  tekanan (pressure drop), profil transfer kalor, serta lima batas
  operasional fisik (batas sonik, batas viskos, batas kapiler, batas
  entrainment, dan batas pendidihan) pada pipa kalor logam alkali suhu
  tinggi?
\item
  Seberapa akurat algoritma komputasi Python yang dikembangkan tersebut
  apabila divalidasi terhadap data batas performa dan titik operasi
  nominal (suhu 1000 K) dari pipa kalor natrium pada purwarupa
  eksperimental \emph{Kilopower Reactor Using Stirling TechnologY}
  (KRUSTY)?
\item
  Bagaimana cara mengaplikasikan model komputasi yang telah tervalidasi
  untuk mengevaluasi kinerja dan mengoptimalkan desain konfigurasi pipa
  kalor pada Reaktor Antariksa Indonesia (ISR)?
\item ~
  \subsection{\texorpdfstring{\textbf{Batasan
  Masalah}}{Batasan Masalah}}\label{batasan-masalah}

  Batasan masalah dari penelitian ini adalah sebagai berikut:
\item
  Pemodelan komputasi perpindahan kalor dan pressure drop dibatasi pada
  kondisi tunak 1-Dimensi (1D \emph{steady-state}) dengan menggunakan
  metode jaringan hambatan termal (\emph{thermal resistance network}).
\item
  Perhitungan komputasi difokuskan untuk menyelesaikan dan memprediksi
  lima batas operasi utama pipa kalor, yaitu: batas sonik, batas viskos,
  batas kapiler, batas \emph{entrainment}, dan batas pendidihan.
\item
  Fluida kerja pipa kalor yang dimodelkan adalah logam alkali, secara
  spesifik menggunakan korelasi sifat termofisik natrium cair
  (\emph{sodium}).
\item
  Validasi model komputasi dibatasi dengan menggunakan titik desain
  fisik dan data eksperimen dasar dari proyek \emph{Kilopower Reactor
  Using Stirling TechnologY} (KRUSTY) pada suhu operasi nominal 1000 K
  (727 °C).
\item
  Aplikasi model ditekankan pada evaluasi konfigurasi pendinginan pasif
  untuk desain konseptual Reaktor Antariksa Indonesia (ISR) yang
  berspektrum cepat dengan daya 500 kWth.
\item
  \textbf{Tujuan Penelitian}

  Tujuan dari penelitian ini adalah sebagai berikut:
\item
  Mengembangkan algoritma pemodelan komputasi berbasis Python yang
  mandiri (\emph{self-contained}) untuk memetakan penurunan tekanan
  (\emph{pressure drop}), transfer kalor, serta lima batas operasional
  fisik pada pipa kalor logam alkali suhu tinggi.
\item
  Memvalidasi keakuratan alat komputasi Python yang dikembangkan dengan
  membandingkannya terhadap batas performa dan profil operasi nominal
  dari reaktor purwarupa KRUSTY.
\item
  Mengaplikasikan model komputasi yang telah tervalidasi untuk
  mengevaluasi kinerja dan mengoptimalkan desain konfigurasi pipa kalor
  pada Reaktor Antariksa Indonesia (ISR).
\item ~
  \subsection{\texorpdfstring{\textbf{Manfaat
  Penelitian}}{Manfaat Penelitian}}\label{manfaat-penelitian}
\item
  Memberikan alternatif alat analisis termal-hidrolik berbasis Python
  yang lebih mudah diakses, fleksibel, dan efisien untuk desain awal
  (preliminer) pipa kalor, mengurangi kebergantungan pada kode komputasi
  berat dan berbiaya tinggi.
\item
  Menghasilkan landasan evaluasi termal-hidrolik yang memastikan
  distribusi panas yang kuat, mitigasi titik panas lokal (hotspots),
  serta memastikan stabilitas aliran kapiler tanpa risiko dryout di
  dalam teras ISR.
\item
  Berkontribusi dalam pengembangan sistem pendinginan pasif yang andal
  dan aman (tanpa komponen bergerak seperti pompa mekanis), yang krusial
  untuk kesuksesan dan umur panjang misi eksplorasi luar angkasa dengan
  reaktor nuklir.
\end{enumerate}

\section{\texorpdfstring{\textbf{BAB II} \textbf{TINJAUAN
PUSTAKA}}{BAB II TINJAUAN PUSTAKA}}\label{bab-ii-tinjauan-pustaka}

Penelitian mengenai aplikasi reaktor nuklir antariksa yang didinginkan
dengan pipa kalor telah mencapai tonggak sejarah penting melalui
keberhasilan eksperimen Kilopower Reactor Using Stirling TechnologY
(KRUSTY) yang diinisiasi oleh NASA dan National Nuclear Security
Administration (NNSA). Proyek purwarupa ini ditujukan untuk menghasilkan
daya listrik sebesar 1 hingga 10 kWe secara kontinu selama belasan tahun
untuk kebutuhan misi luar angkasa. Pada uji coba daya penuh di tahun
2018, KRUSTY berhasil beroperasi dengan memanfaatkan teras padat
Uranium-235 yang memindahkan panas fisi sebesar kurang lebih 5,5 kWth
menuju mesin konversi Stirling menggunakan jaringan pipa kalor natrium
pasif. Keberhasilan eksperimen KRUSTY ini menjadikannya standar
pembuktian (benchmark) yang paling mutakhir dan ideal untuk memvalidasi
pemodelan reaktor antariksa, mengingat kemampuannya bertahan dalam
berbagai kondisi kegagalan simulasi dan transien operasi

Sementara itu, dari sisi pengembangan reaktor nasional, Khandaq dkk.
telah mengusulkan studi desain konseptual teras Indonesian Space Reactor
(ISR). ISR dirancang sebagai reaktor berspektrum neutron cepat dengan
target keluaran panas sebesar 500 kWth untuk masa operasi lebih dari 10
tahun. Teras reaktor ini berstruktur silinder heksagonal berongga yang
terdiri atas 61 pin bahan bakar uranium nitrat (UN) dengan pengayaan
U-235 sekitar 55\%. Pada gagasan aslinya, pin bahan bakar ISR didesain
dalam bentuk anulus di mana fluida pendingin logam cair (Na-78K)
mengalir secara aktif di dalamnya untuk meratakan distribusi panas
radial. Literatur ini sangat esensial bagi penelitian skripsi ini karena
memberikan referensi spesifikasi geometri dan beban termal yang akan
dikonversi dan dievaluasi kinerjanya menggunakan konsep pendinginan
pasif pipa kalor

Dalam bidang pemodelan komputasional pipa kalor itu sendiri, terdapat
berbagai perangkat lunak khusus yang telah dijadikan standar industri.
Secara historis, desainer termal-hidrolik menggunakan kode lawas seperti
ANLHTP buatan Argonne National Laboratory atau kode HTPIPE berbasis
Fortran dari Los Alamos National Laboratory (LANL) yang difokuskan untuk
memprediksi profil suhu dan menyelesaikan batas operasional keadaan
tunak (\emph{steady-state}) secara 1-Dimensi. Pada perkembangan modern,
perangkat lunak dengan komputasi masif seperti Sockeye yang dibangun di
atas kerangka \emph{Multiphysics Object-Oriented Simulation Environment}
(MOOSE) digunakan untuk simulasi transien yang menggabungkan aliran uap
1D dan konduksi termal padatan 2D secara mekanistik. Keterbatasan
fleksibilitas pada kode lawas serta mahalnya beban komputasi pada
Sockeye mengindikasikan bahwa penyusunan skrip analitik yang mandiri
(\emph{self-contained}) dan efisien menggunakan bahasa pemrograman
Python sangat dibutuhkan untuk analisis desain awal

Sebagai alternatif pengembangan instrumen analisis, studi terkini yang
mengandalkan bahasa pemrograman tingkat tinggi untuk pemodelan reaktor
KRUSTY telah dilakukan oleh Zhang dkk. dengan menggunakan perangkat
MATLAB/Simulink. Mereka mengembangkan lingkungan pemodelan multi-domain
yang mengintegrasikan dinamika reaktor titik dan modul perpindahan panas
pipa kalor untuk mengevaluasi operasi transien reaktor. Validasi alat
komputasi tersebut terhadap data eksperimen KRUSTY menunjukkan galat
suhu yang berada di bawah 3\% pada sebagian besar tahapan operasi,
membuktikan keandalan pemrograman matematika untuk simulasi reaktor.
Senada dengan hal tersebut, Guo dkk. juga mensimulasikan distribusi suhu
teras KRUSTY melalui pengembangan \emph{Super Thermal Conductivity
Model} (STCM) di mana transportasi uap di dalam pipa disederhanakan
menjadi model jaringan hambatan termal.

Keseluruhan literatur di atas mengukuhkan landasan bagi penelitian
skripsi ini. Eksperimen KRUSTY menyediakan landasan data empiris yang
terbukti sukses, sementara desain ISR menawarkan kasus aplikasi reaktor
antariksa tingkat lanjut dengan daya lebih besar yang menuntut
perancangan termal-hidrolik yang presisi. Dengan mengidentifikasi
kesenjangan ketersediaan alat analisis pada kode-kode terdahulu seperti
HTPIPE maupun perangkat berat seperti Sockeye, penelitian ini
diposisikan untuk membangun model komputasi berbasis Python yang
mengintegrasikan penyelesaian batas operasional dan evaluasi perpindahan
kalor. Melalui komparasi dengan performa KRUSTY, model Python yang
dihasilkan akan memiliki justifikasi yang valid sebelum diekstrapolasi
untuk menguji kelayakan integrasi pipa kalor pada desain Reaktor
Antariksa Indonesia (ISR).

\section{\texorpdfstring{\textbf{BAB III} \textbf{DASAR
TEORI}}{BAB III DASAR TEORI}}\label{bab-iii-dasar-teori}

\begin{enumerate}
\def\labelenumi{\arabic{enumi}.}
\item ~
  \subsection{\texorpdfstring{\textbf{Prinsip Kesetimbangan Tekanan Pipa
  Kalor}}{Prinsip Kesetimbangan Tekanan Pipa Kalor}}\label{prinsip-kesetimbangan-tekanan-pipa-kalor}
\end{enumerate}

Untuk mempertahankan transportasi panas pada kondisi tunak
(\emph{steady-state}), sebuah pipa kalor harus memenuhi kriteria
kesetimbangan tekanan dinamis global. Proses penguapan yang terjadi di
bagian evaporator menghasilkan gradien tekanan yang mendorong uap
mengalir secara aksial menuju kondensor. Di bagian kondensor, uap
melepaskan panas latennya, berubah fase menjadi cairan, dan ditarik
kembali menuju evaporator oleh gaya pompa kapiler (\emph{capillary
pumping head}) yang dihasilkan oleh struktur sumbu berpori (\emph{wick})

Selama sirkulasi tersebut, fluida kerja mengalami gaya gesek
(\emph{frictional forces}) baik pada fase uap maupun cairan. Agar
sirkulasi fluida tidak terputus dan fenomena kekeringan (\emph{dryout})
pada evaporator dapat dihindari, tekanan kapiler maksimum yang mampu
dihasilkan oleh sumbu (\(\Delta {P}_{max,cap}\)) harus lebih besar atau
sama dengan total penurunan tekanan di seluruh sistem:

\href{https://www.codecogs.com/eqnedit.php?latex=\%5CDelta\%7B\%7DP_\%7Bcap\%2Cmax\%7D\%5Cgeq\%7B\%7D\%5CDelta\%7B\%7DP_\%7Bv\%7D\%2B\%5CDelta\%7B\%7DP_\%7Bl\%7D\%2B\%5CDelta\%7B\%7DP_\%7Bg\%7D\#0}{}

Dengan \(\Delta {P}_{v}\) adalah penurunan tekanan aliran uap,
\(\Delta {P}_{l}\) adalah penurunan tekanan cairan di dalam sumbu, dan
\(\Delta {P}_{g}\) adalah kerugian tekanan hidrostatik akibat gravitasi.

Gaya pompa kapiler maksimum itu sendiri dievaluasi menggunakan persamaan
Young yang bergantung pada tegangan permukaan fluida \(\sigma\), sudut
kontak \(\theta\), dan radius pori efektif sumbu \({r}_{c}\):

\href{https://www.codecogs.com/eqnedit.php?latex=\%5CDelta\%7B\%7DP_\%7Bcap\%2Cmax\%7D\%3D\%5Cfrac\%7B2\%5Csigma\%7B\%7D\%5Ccos\%7B\%5Ctheta\%7B\%7D\%7D\%7D\%7Br_\%7Bc\%7D\%7D\#0}{}

\begin{enumerate}
\def\labelenumi{\arabic{enumi}.}
\setcounter{enumi}{1}
\item ~
  \subsection{\texorpdfstring{\textbf{Pemodelan Penurunan
  Tekanan}}{Pemodelan Penurunan Tekanan}}\label{pemodelan-penurunan-tekanan}
\end{enumerate}

Aliran cairan yang kembali melalui struktur sumbu berpori dimodelkan
sebagai aliran laminar. Dengan memanfaatkan Hukum Darcy untuk media
berpori yang diintegrasikan sepanjang pipa kalor, total penurunan
tekanan cairan, \(\Delta {P}_{l}\) dapat diekspresikan sebagai fungsi
dari laju perpindahan kalor (\(Q\)) menggunakan hukum Darcy:

\href{https://saxarona.github.io/mathjax-viewer/?input=\%5CDelta\%7B\%7DP_\%7Bl\%7D\%3D\%5Cfrac\%7B\%5Cmu_\%7Bl\%7Dl_\%7Beff\%7DQ\%7D\%7B\%5Crho\%7B\%7D_\%7Bl\%7DKA_\%7Bw\%7Dh_\%7Bfg\%7D\%7D\#0}{}

Di mana

\begin{itemize}
\item
  \emph{\({\mu }_{l}\)} adalah viskositas dinamis cairan (Pa·s)
\item
  \({\rho }_{l}\) adalah densitas cairan (kg/m3)
\item
  \({h}_{fg}\) adalah kalor laten penguapan (J/kg)
\item
  \({A}_{w}\) adalah luas penampang melintang struktur sumbu (m2)
\item
  \(K\) merupakan permeabilitas hidrolik dari sumbu tersebut (m2)
\end{itemize}

Variabel \({l}_{eff}\) adalah panjang efektif pipa kalor yang
memperhitungkan distribusi spasial dari penguapan dan kondensasi,
dirumuskan dengan:

\href{https://www.codecogs.com/eqnedit.php?latex=l_\%7Beff\%7D\%3D0.5l_e\%2Bl_a\%2B0.5l_c\#0}{}

Di mana \({l}_{e}\), \({l}_{a}\), dan \({l}_{c}\) secara berturut-turut
merupakan panjang area evaporator, adiabatik, dan kondensor.

\begin{enumerate}
\def\labelenumi{\arabic{enumi}.}
\setcounter{enumi}{2}
\item ~
  \subsection{\texorpdfstring{\textbf{Batasan Kerja Termal-Hidrolik
  (\emph{Operational
  Limits})}}{Batasan Kerja Termal-Hidrolik (Operational Limits)}}\label{batasan-kerja-termal-hidrolik-operational-limits}
\end{enumerate}

Dalam pengoperasian pipa kalor, laju perpindahan kalor maksimum yang
dapat dicapai dibatasi oleh interaksi dinamika fluida dan termodinamika
fase cair-uap di dalamnya.

Karakteristik limitasi ini bervariasi bergantung pada geometri pipa,
sifat termofisika fluida kerja, struktur sumbu kapiler (wick), serta
temperatur operasi. Kegagalan operasional atau penurunan drastis laju
perpindahan kalor terjadi apabila salah satu batas hidrodinamika atau
termal terlampaui. Secara umum, limitasi kerja termal-hidrolik pada pipa
kalor mencakup batasan viskos, sonik, \emph{entrainment}, kapiler, dan
pendidihan, yang masing-masing dijabarkan sebagai berikut.

\begin{enumerate}
\def\labelenumi{\arabic{enumi}.}
\item ~
  \subsubsection{\texorpdfstring{\textbf{Batasan
  Viskos}}{Batasan Viskos}}\label{batasan-viskos}
\end{enumerate}

Batas viskos sangat relevan ketika reaktor antariksa baru mulai
dihidupkan (startup) dari keadaan dingin atau beku. Pada suhu operasi
yang sangat rendah, tekanan uap berada pada kondisi yang sangat kecil
sehingga gaya viskos (gesekan internal fluida) menjadi dominan. Batas
viskos terjadi apabila gaya gesek viskos uap mengalahkan gradien tekanan
yang dihasilkan oleh proses penguapan, sehingga uap mengalami stagnasi
dan aliran sirkulasi tidak dapat berlangsung. Dengan mengasumsikan uap
gas ideal pada aliran satu dimensi, aliran inkompresibel, serta beban
penguapan dan kondensasi yang seragam, batas daya viskos (\({Q}_{vis}\))
dapat diprediksi melalui persamaan:

\href{https://saxarona.github.io/mathjax-viewer/?input=Q_\%7Bvis\%7D\%3D\%5Cfrac\%7BA_\%7Bv\%7Dr_\%7Bv\%7D\%5E2h_\%7Bfg\%7D\%5Crho_\%7Bv\%7DP_\%7Bv\%7D\%7D\%7B16\%5Cmu_\%7Bv\%7Dl_\%7Beff\%7D\%7D\#0}{}

di mana

\begin{itemize}
\item
  \({A}_{v}\) adalah luas penampang inti uap (m2)\\
\item
  \({r}_{v}\) adalah radius ruang uap (m)\\
\item
  \({h}_{fg}\) adalah kalor laten penguapan (J/kg)\\
\item
  \({\rho }_{v}\) adalah densitas uap saturasi (kg/m3)\\
\item
  \({P}_{v}\) adalah tekanan uap saturasi (Pa)\\
\item
  \({\mu }_{v}\) adalah viskositas dinamis uap (Pa·s)\\
\item
  \({l}_{eff}\) adalah panjang efektif pipa kalor (m)

  2. \#\#\# \textbf{Batasan Sonik}
\end{itemize}

Batas sonik tercapai apabila kecepatan aliran uap di ujung pintu keluar
evaporator menyamai kecepatan suara lokal (\(M=1\)), sehingga
menyebabkan terjadinya fenomena aliran tercekik (\emph{choked flow}).
Berdasarkan model aliran inersia satu dimensi, kapasitas perpindahan
kalor tercekik maksimum (\({Q}_{sonic}\)) dirumuskan sebagai:

\href{https://saxarona.github.io/mathjax-viewer/?input=Q_\%7Bsonic\%7D\%3DA_v\%5Crho\%7B\%7D_vh_\%7Bfg\%7D\%5Cleft\%5B\%5Cfrac\%7B\%5Cgamma\%7B\%7DR_vT_v\%7D\%7B2(\%5Cgamma\%2B1)\%7D\%5Cright\%5D\%5E\%7B0.5\%7D\#0}{}

di mana

\begin{itemize}
\tightlist
\item
  \({A}_{v}\) adalah luas penampang inti uap (m²)\\
\item
  \({\rho }_{v}\) adalah densitas uap saturasi (kg/m³)\\
\item
  \({h}_{fg}\) adalah kalor laten penguapan (J/kg)\\
\item
  \(\gamma\) adalah rasio kapasitas kalor (tanpa satuan,
  \({c}_{p}/{c}_{v}\))\\
\item
  \({R}_{v}\) adalah konstanta gas spesifik uap (J/(kg·K))\\
\item
  \({T}_{v}\) adalah suhu uap saturasi evaporator (K)
\end{itemize}

Melampaui batas ini tidak selalu menyebabkan kerusakan permanen pada
pipa kalor seperti halnya batas pendidihan, melainkan akan membatasi
kemampuan efisiensi perpindahan kalor hingga suhu evaporator meningkat
dan tekanan uap naik cukup tinggi untuk keluar dari rezim sonik.

\begin{enumerate}
\def\labelenumi{\arabic{enumi}.}
\setcounter{enumi}{2}
\item ~
  \subsubsection{\texorpdfstring{\textbf{Batasan
  Entrainment}}{Batasan Entrainment}}\label{batasan-entrainment}
\end{enumerate}

Di dalam pipa kalor, cairan mengalir melalui struktur sumbu kapiler
secara berlawanan arah terhadap aliran uap berkecepatan tinggi pada inti
ruang uap. Interaksi antara kedua fase aliran ini menimbulkan gaya geser
yang signifikan pada antarmuka cair-uap. Batasan \emph{entrainment}
tercapai apabila gaya geser uap melampaui gaya tegangan permukaan
cairan, sehingga butiran cairan terlepas dari struktur sumbu, ikut
terbawa oleh aliran uap, dan terakumulasi di area kondensor. Fenomena
ini dianalisis menggunakan rasio antara gaya inersia uap terhadap gaya
tegangan permukaan yang dinyatakan melalui bilangan Weber (Weber
number). Kapasitas perpindahan kalor maksimum akibat fenomena
entrainment (\({Q}_{e}\)) dirumuskan sebagai berikut:

\href{https://saxarona.github.io/mathjax-viewer/?input=Q_e\%3DA_vh_\%7Bfg\%7D\%5Cleft\%5B\%5Cfrac\%7B\%5Csigma\%7B\%7D\%5Crho_v\%7D\%7B2r_h\%7D\%5Cright\%5D\%5E\%7B0.5\%7D\#0}{}

dengan

\begin{itemize}
\item
  \({A}_{v}\) adalah luas penampang ruang uap (m²)\\
\item
  \({h}_{fg}\) adalah kalor laten penguapan (J/kg)\\
\item
  \(\sigma\) adalah tegangan permukaan cairan (N/m)\\
\item
  \({\rho }_{v}\) adalah densitas uap saturasi (kg/m³)\\
\item
  \({r}_{h}\) adalah radius hidrolik \emph{wick} sumbu kapiler (m)

  4. \#\#\# \textbf{Batasan Pendidihan}

  Berbeda dengan batasan operasional lainnya yang membatasi fluks panas
  aksial, batas pendidihan membatasi fluks panas secara radial pada
  evaporator. Jika fluks panas radial yang diterima dinding evaporator
  terlalu ekstrem, gradien suhu konduksi melintasi struktur sumbu dan
  cairan akan memicu kondisi \emph{superheat}. Hal ini menyebabkan
  terjadinya pendidihan nukleat di dalam struktur sumbu, sehingga
  gelembung uap yang terbentuk akan menghambat pasokan balik cairan
  kapiler dan memicu dryout lokal serta lonjakan suhu atau
  \emph{hotspot} pada teras reaktor. Kapasitas perpindahan kalor
  maksimum untuk mencegah fenomena batas pendidihan (\({Q}_{b}\))
  ditentukan berdasarkan pendekatan nukleasi gelembung uap:
\end{itemize}

\href{https://saxarona.github.io/mathjax-viewer/?input=Q_b\%3D\%5Cfrac\%7B2\%5Cpi\%7B\%7Dl_ek_\%7Beff\%7DT_v\%7D\%7Bh_\%7Bfg\%7D\%5Crho_\%7Bv\%7D\%5Cln(r_i\%2Fr_v)\%7D\%5Cleft\%5B\%5Cfrac\%7B2\%5Csigma\%7B\%7D\%7D\%7Br_n\%7D-P_c\%5Cright\%5D\#0}{}

di mana

\begin{itemize}
\tightlist
\item
  \({l}_{e}\) adalah panjang evaporator (m)\\
\item
  \({k}_{eff}\) adalah konduktivitas termal efektif sumbu (W/(m·K))\\
\item
  \({T}_{v}\) adalah suhu uap (K)\\
\item
  \(\sigma\) adalah tegangan permukaan (N/m)\\
\item
  \({h}_{fg}\) adalah kalor laten penguapan (J/kg)\\
\item
  \({\rho }_{v}\) adalah densitas uap (kg/m³)\\
\item
  \({r}_{i}\) adalah radius dalam pipa (m)\\
\item
  \({r}_{v}\) adalah radius inti uap (m)\\
\item
  \({r}_{n}\) adalah radius nukleasi gelembung kritis (m)\\
\item
  \({P}_{c}\) adalah tekanan kapiler lokal (Pa)\\
\end{itemize}

\begin{enumerate}
\def\labelenumi{\arabic{enumi}.}
\setcounter{enumi}{3}
\tightlist
\item
  \textbf{Sifat Termofisik Natrium Cair}
\end{enumerate}

Mengingat semua perhitungan kerugian tekanan dan batas operasional
sangat bergantung pada kondisi fluida, pemodelan komputasi memerlukan
persamaan termofisik yang akurat. Sifat termofisik natrium jenuh
berwujud sangat non-linear terhadap perubahan suhu operasi. Berdasarkan
korelasi empiris dari \emph{Argonne National Laboratory}, densitas
cairan natrium (\({\rho }_{l}\)) dalam kg/m3 dari rentang titik lelehnya
hingga 1644 K dihitung melalui polinomial:

\href{https://www.codecogs.com/eqnedit.php?latex=\%5Crho\%7B\%7D_l\%3D1011.8-0.22054T-1.9226\%5Ctimes\%7B\%7D10\%5E\%7B-5\%7DT\%5E2\%2B5.6371\%5Ctimes\%7B\%7D10\%5E\%7B-9\%7DT\%5E3\#0}{}

Sementara itu, tegangan permukaan (\(\sigma\)) dalam satuan N/m, yang
menjadi parameter penentu daya pompa kapiler, dimodelkan menggunakan
persamaan linear

\begin{enumerate}
\def\labelenumi{\arabic{enumi}.}
\setcounter{enumi}{4}
\tightlist
\item
  \textbf{Analisis Jaringan Hambatan Termal}
\end{enumerate}

Untuk mengevaluasi profil penurunan suhu di sepanjang pipa kalor (dari
luar dinding evaporator hingga luar dinding kondensor), penelitian ini
mengadopsi pemodelan jaringan hambatan termal (thermal resistance
network). Pendekatan ini menganalogikan aliran panas dengan sirkuit
listrik. Total resistansi termal merupakan kombinasi perhitungan
hambatan secara seri dan paralel, yang secara umum terdiri atas sembilan
hambatan: hambatan konduksi arah radial di dinding dan sumbu pada area
evaporator, batas antarmuka transisi fase cair-uap, hambatan aliran uap
di zona adiabatik, serta resistansi arah radial di sumbu dan dinding
kondensor. Dengan menyusun rangkaian ini dalam pemrograman Python,
gradien suhu aktual untuk fluks beban dari reaktor dapat disimulasikan

\textbf{BAB IV}\\
\textbf{PELAKSANAAN PENELITIAN}

LOREM IPSUM

\textbf{BAB V}\\
\textbf{HASIL DAN PEMBAHASAN}

LOREM IPSUM

\textbf{BAB VI}\\
\textbf{KESIMPULAN DAN SARAN}

\begin{enumerate}
\def\labelenumi{\arabic{enumi}.}
\tightlist
\item
  \textbf{Kesimpulan}
\end{enumerate}

LOREM IPSUM

\begin{enumerate}
\def\labelenumi{\arabic{enumi}.}
\setcounter{enumi}{1}
\tightlist
\item
  \textbf{Saran}
\end{enumerate}

LOREM IPSUM

\textbf{DAFTAR PUSTAKA}

\href{https://www.zotero.org/google-docs/?cab6zL}{{[}1{]}}
\href{https://www.zotero.org/google-docs/?cab6zL}{P. R. McClure, D. I.
Poston, M. A. Gibson, L. S. Mason, and R. C. Robinson, `Kilopower
Project: The KRUSTY Fission Power Experiment and Potential Missions',
\emph{Nucl. Technol.}, vol.~206, no. sup1, pp.~S1--S12, Jun.~2020, doi:
10.1080/00295450.2020.1722554.}

\href{https://www.zotero.org/google-docs/?cab6zL}{{[}2{]}}
\href{https://www.zotero.org/google-docs/?cab6zL}{D. Beard, C. Tarau,
and W. G. Anderson, `Sodium Heat Pipes for Space and Surface Fission
Power', in \emph{15th International Energy Conversion Engineering
Conference}, Atlanta, GA: American Institute of Aeronautics and
Astronautics, Jul.~2017. doi: 10.2514/6.2017-5087.}

\href{https://www.zotero.org/google-docs/?cab6zL}{{[}3{]}}
\href{https://www.zotero.org/google-docs/?cab6zL}{H. Jouhara, D. A.
Reay, R. J. McGlen, P. Kew, and J. McDonough, \emph{Heat Pipes: theory,
design and applications}, Seventh edition. Oxford Cambridge, MA:
Butterworth-Heinemann, an imprint of Elsevier, 2024.}

\end{document}
